\section{Trabajo futuro y evaluación pendiente}
\label{sec:trabajo-futuro}
La continuación del trabajo se concentra en evaluar la segunda etapa bajo condiciones documentadas. Los registros de SFT45k y del piloto permiten analizar el aprendizaje descrito en el capítulo~\ref{ch:evolucion-entrenamiento}, mientras que la evaluación externa de la segunda etapa y el contraste con RaR requieren completarse. Por ello, las rúbricas preparadas y las configuraciones implementadas deben distinguirse de los resultados que todavía no se presentan.

En primer lugar, corresponde comprobar cómo los criterios construidos a partir de las obligaciones de la consulta se comportan ante respuestas correctas, parciales, contradictorias y equivalentes en significado. Esta revisión debe permitir calibrar al juez y detectar recompensas que favorezcan menciones superficiales o exigencias no solicitadas antes de interpretar una mayor puntuación como una mejora de las respuestas.

A continuación, las variantes de recompensa y de ayudas deberán compararse dentro de la segunda etapa, manteniendo constantes el modelo de partida, los datos y la evaluación. Estas comparaciones permitirán estudiar por separado el efecto de las rúbricas, de su puntuación y de las ayudas. El contraste con RaR también requiere comprobar las condiciones descritas en la sección~\ref{sec:comparacion-rar} y explicitar aquellas que no sean equivalentes.

Por último, el experimento final previsto en la sección~\ref{sec:protocolo-final-previsto} contempla un reparto del conjunto de aprendizaje en mitades para SFT y GRPO, con validación y prueba reservadas. Esta propuesta permitirá evaluar el procedimiento bajo una separación definida desde el inicio, pero no se considera ejecutada ni se confunde con el SFT45k o con el piloto preparado. Las curvas de entrenamiento y los resultados externos deberán incorporarse cuando estén disponibles y verificados.
